\documentclass[aspectratio=169,11pt]{beamer}

\usetheme{default}
\setbeamertemplate{navigation symbols}{}
\setbeamertemplate{footline}[frame number]

\usepackage{amsmath}
\usepackage{amsfonts}
\usepackage{booktabs}
\usepackage{graphicx}
\graphicspath{{master/figures/}{figures/}}

\definecolor{hcblue}{RGB}{30,78,121}
\definecolor{hcgold}{RGB}{193,148,37}
\definecolor{hcslate}{RGB}{59,70,84}
\definecolor{hclight}{RGB}{240,244,248}

\setbeamercolor{title}{fg=white,bg=hcblue}
\setbeamercolor{frametitle}{fg=white,bg=hcblue}
\setbeamercolor{structure}{fg=hcblue}
\setbeamercolor{block title}{fg=white,bg=hcblue}
\setbeamercolor{block body}{fg=black,bg=hclight}
\setbeamercolor{normal text}{fg=black,bg=white}
\setbeamercolor{itemize item}{fg=hcgold}
\setbeamercolor{itemize subitem}{fg=hcgold}
\setbeamerfont{title}{series=\bfseries,size=\LARGE}
\setbeamerfont{frametitle}{series=\bfseries,size=\large}

\title{Hierarchical Clustering}
\subtitle{Intuition, dendrograms, linkage, and the mathematics of greedy merges}
\author{}
% \date{March 31, 2026}

\newcommand{\ESS}{\mathrm{ESS}}

\begin{document}

\begin{frame}[plain]
  \titlepage
\end{frame}

\begin{frame}[t]{Why Hierarchical Clustering?}
  \small
  \begin{columns}[T,totalwidth=\textwidth]
    \column{0.48\textwidth}
      \begin{itemize}
        \item Hierarchical clustering is \textbf{unsupervised} and \textbf{connectivity-based}.
        \item It builds a \textbf{nested family of clusters}, not just one partition.
        \item We do \emph{not} need to pick one value of $K$ up front.
        \item The raw ingredients are observations plus a dissimilarity:
        \[
          \text{observations} \;+\; \text{a dissimilarity } d(x_i,x_j).
        \]
        \item Nearby points merge low in the tree; distant groups only merge high up.
      \end{itemize}
    \column{0.48\textwidth}
      \begin{center}
        \includegraphics[
          width=\linewidth,
          height=0.52\textheight,
          keepaspectratio
        ]{agglomerative_vs_divisive.jpeg}
      \end{center}
  \end{columns}
\end{frame}

\begin{frame}[t]{Agglomerative Versus Divisive}
  \begin{columns}[T,totalwidth=\textwidth]
    \column{0.47\textwidth}
      \begin{block}{Agglomerative}
        \begin{itemize}
          \item Bottom-up
          \item Start with $n$ singleton clusters
          \item Repeatedly merge the two closest clusters
          \item Most common form in practice
        \end{itemize}
      \end{block}

      \[
        \text{singletons}
        \;\longrightarrow\;
        \text{larger and larger groups}
      \]
    \column{0.47\textwidth}
      \begin{block}{Divisive}
        \begin{itemize}
          \item Top-down
          \item Start with one large cluster containing every point
          \item Repeatedly split a cluster into smaller pieces
          \item Conceptually natural, but less common in standard software
        \end{itemize}
      \end{block}

      \[
        \text{one big cluster}
        \;\longrightarrow\;
        \text{finer and finer groups}
      \]
  \end{columns}

  \vspace{0.35em}
  \begin{itemize}
    \item In this deck, ``hierarchical clustering'' means \textbf{agglomerative hierarchical clustering} unless stated otherwise.
  \end{itemize}
\end{frame}

\begin{frame}[t]{How to Read a Dendrogram}
  \begin{columns}[T,totalwidth=\textwidth]
    \column{0.50\textwidth}
      \begin{center}
        \includegraphics[
          width=\linewidth,
          height=0.63\textheight,
          keepaspectratio
        ]{dendrogram_parts.png}

        \vspace{0.25em}
        \scriptsize Adapted from \texttt{refs/10-hierarchical-clustering.pptx}
      \end{center}
    \column{0.46\textwidth}
      \begin{itemize}
        \item Each \textbf{leaf} is one observation.
        \item A horizontal \textbf{link} marks the moment when two clusters merge.
        \item The \textbf{height} of that link is the dissimilarity at which the merge happened.
        \item Clusters that join low in the tree are more similar than clusters that join high in the tree.
      \end{itemize}

      \begin{block}{Interpretation}
        The vertical axis matters much more than the horizontal ordering of the leaves.
      \end{block}
  \end{columns}
\end{frame}

\begin{frame}[t]{Interpreting Dendrograms Carefully}
  \begin{columns}[T,totalwidth=\textwidth]
    \column{0.52\textwidth}
      \begin{center}
        \includegraphics[
          width=\linewidth,
          height=0.52\textheight,
          keepaspectratio
        ]{islp_dendrogram_interpretation.png}

        \vspace{0.25em}
        \scriptsize Adapted from \texttt{quizzes/book-snips/ISLP-12.4-Clustering.pdf}
      \end{center}
    \column{0.44\textwidth}
      \small
      \begin{itemize}
        \item The \textbf{height} of a merge is meaningful; the left-to-right ordering of leaves usually is not.
        \item A point may sit horizontally near another point in the dendrogram but still merge with it at a much larger height.
        \item Always connect the dendrogram back to the original feature space or distance matrix before making a story out of it.
      \end{itemize}

      \begin{block}{Rule of thumb}
        Read dendrograms vertically, not horizontally.
      \end{block}
  \end{columns}
\end{frame}

\begin{frame}[t]{Agglomerative Clustering Algorithm}
  \begin{columns}[T,totalwidth=\textwidth]
    \column{0.48\textwidth}
      \scriptsize
      Start with a pairwise dissimilarity matrix
      \[
        D^{(0)}_{ij}=d(x_i,x_j), \qquad D^{(0)}_{ii}=0.
      \]

      \begin{block}{Greedy merge rule}
        At step $t$, choose
        \[
          (A_t,B_t)
          =
          \arg\min_{A\neq B} d_{\text{link}}(A,B),
        \]
        merge them, and record the merge height
        \[
          h_t = d_{\text{link}}(A_t,B_t).
        \]
      \end{block}

      Initialize with $n$ singleton clusters, repeatedly merge and update, and stop when one cluster remains.

      \vspace{0.4em}
      The method is \textbf{greedy}: once two clusters merge, that choice cannot be undone.
    \column{0.48\textwidth}
      \begin{center}
        \includegraphics[
          width=\linewidth,
          height=0.54\textheight,
          keepaspectratio
        ]{hc_first_steps_complete_linkage.png}
      \end{center}
  \end{columns}
\end{frame}

\begin{frame}[t]{The Proximity Matrix Drives the Merges}
  \begin{columns}[T,totalwidth=\textwidth]
    \column{0.50\textwidth}
      \footnotesize
      For observations $x_1,\ldots,x_n$, start with the pairwise dissimilarity matrix
      \[
        D=(d_{ij})_{i,j=1}^n,
        \qquad
        d_{ij}=d(x_i,x_j).
      \]

      \[
      D=
      \begin{bmatrix}
      0 & d_{12} & \cdots & d_{1n}\\
      d_{21} & 0 & \cdots & d_{2n}\\
      \vdots & \vdots & \ddots & \vdots\\
      d_{n1} & d_{n2} & \cdots & 0
      \end{bmatrix}
      \]

      \begin{itemize}
        \item It is symmetric when $d(x_i,x_j)=d(x_j,x_i)$.
        \item Diagonal entries are $0$ because each point has zero distance from itself.
      \end{itemize}
    \column{0.45\textwidth}
      \footnotesize
      \begin{block}{What changes after a merge?}
        If $A$ and $B$ merge into $A\cup B$, then every remaining cluster $C$ needs a new distance
        \[
          d_{\text{link}}(A\cup B, C).
        \]
      \end{block}

      \begin{itemize}
        \item After each merge, two rows/columns are replaced by one for the merged cluster.
        \item Single linkage uses the closest cross-cluster pair.
        \item Complete linkage uses the farthest cross-cluster pair.
        \item Average linkage uses the mean cross-cluster distance.
      \end{itemize}
  \end{columns}
\end{frame}

\begin{frame}[t]{The Key Mathematical Choice: Linkage}
  \begin{columns}[T,totalwidth=\textwidth]
    \column{0.52\textwidth}
      \small
      For clusters $A$ and $B$, the linkage function defines their distance:

      \[
        d_{\text{single}}(A,B)
        =
        \min_{i\in A,\;j\in B} d(x_i,x_j)
      \]
      \[
        d_{\text{complete}}(A,B)
        =
        \max_{i\in A,\;j\in B} d(x_i,x_j)
      \]
      \[
        d_{\text{average}}(A,B)
        =
        \frac{1}{|A||B|}
        \sum_{i\in A}\sum_{j\in B} d(x_i,x_j)
      \]
      \[
        d_{\text{centroid}}(A,B)
        =
        \|\bar x_A-\bar x_B\|_2
      \]
      \[
        d_{\text{Ward}}(A,B)
        =
        \ESS(A\cup B)-\ESS(A)-\ESS(B)
      \]
    \column{0.44\textwidth}
      \begin{center}
        \includegraphics[
          width=\linewidth,
          height=0.72\textheight,
          keepaspectratio
        ]{linkage_methods_overview.png}

        \vspace{0.25em}
        \scriptsize Adapted from \texttt{refs/10-hierarchical-clustering.pptx}
      \end{center}
  \end{columns}
\end{frame}

\begin{frame}[t]{Ward's Linkage: The Intuition}
  \begin{columns}[T,totalwidth=\textwidth]
    \column{0.54\textwidth}
      \small
      Ward's method does \textbf{not} ask which two clusters have the smallest point-to-point distance.
      Instead, it asks:

      \vspace{0.4em}
      \begin{block}{Guiding question}
        If we merge clusters $A$ and $B$, how much more \textbf{spread out} will the data become?
      \end{block}

      \begin{itemize}
        \item A good merge should keep the new cluster \textbf{tight and compact}.
        \item A bad merge would glue together groups that create a large, diffuse cloud.
        \item So Ward picks the pair whose merge causes the \textbf{smallest increase} in within-cluster variation.
      \end{itemize}

      \[
        \Delta(A,B)=\ESS(A\cup B)-\ESS(A)-\ESS(B)
      \]
    \column{0.40\textwidth}
      \scriptsize
      \begin{block}{Compared with other linkage rules}
        \begin{itemize}
          \item \textbf{Single}: nearest pair.
          \item \textbf{Complete}: farthest pair.
          \item \textbf{Average}: mean pairwise distance.
          \item \textbf{Ward}: the \textbf{whole cluster geometry}.
        \end{itemize}
      \end{block}

      \begin{block}{Practical effect}
        Ward often avoids chaining and tends to prefer rounder, more balanced clusters.
      \end{block}
  \end{columns}
\end{frame}

\begin{frame}[t]{A Tiny Worked Example of Linkage Values}
  \begin{columns}[T,totalwidth=\textwidth]
    \column{0.47\textwidth}
      \footnotesize
      Suppose we have two clusters on the real line:
      \[
        A=\{0,2\}, \qquad B=\{5,6\}.
      \]
      The cross-cluster distances are
      \[
        5,\;6,\;3,\;4.
      \]

      \begin{block}{Linkage values}
        \[
          d_{\text{single}}(A,B)=3
        \]
        \[
          d_{\text{complete}}(A,B)=6
        \]
        \[
          d_{\text{average}}(A,B)=\frac{5+6+3+4}{4}=4.5
        \]
        \[
          d_{\text{centroid}}(A,B)=|1-5.5|=4.5
        \]
      \end{block}
    \column{0.49\textwidth}
      \footnotesize
      For Ward linkage,
      \[
        \bar x_A = 1, \qquad \bar x_B = 5.5.
      \]
      Using the Euclidean simplification,
      \[
        \Delta(A,B)
        =
        \frac{|A||B|}{|A|+|B|}
        \|\bar x_A-\bar x_B\|_2^2
        =
        \frac{2\cdot 2}{4}(4.5)^2
        =
        20.25.
      \]

      \begin{block}{Takeaway}
        The \emph{same two clusters} can look close or far apart depending on the linkage rule, which is why linkage choice strongly affects the dendrogram.
      \end{block}
  \end{columns}
\end{frame}

\begin{frame}[t]{Distance Choice Matters Too}
  \begin{columns}[T,totalwidth=\textwidth]
    \column{0.47\textwidth}
      \begin{block}{If scales differ, standardize first}
        \[
          z_{ij}
          =
          \frac{x_{ij}-\bar x_j}{s_j}
        \]
      \end{block}

      \begin{itemize}
        \item Otherwise, variables with large units can dominate the distance calculation.
        \item A spending variable measured in dollars can overwhelm an age variable measured in years.
        \item Standardization makes each feature contribute on a comparable scale.
      \end{itemize}
    \column{0.47\textwidth}
      \begin{block}{Pick the dissimilarity to match the question}
        \[
          d_{\text{Euclidean}}(x_i,x_j)
          =
          \|x_i-x_j\|_2
        \]
        \[
          d_{\text{corr}}(x_i,x_j)
          =
          1-\mathrm{corr}(x_i,x_j)
        \]
      \end{block}

      \begin{itemize}
        \item Euclidean distance emphasizes absolute magnitude.
        \item Correlation-based distance emphasizes \textbf{shape} or profile similarity.
        \item The right distance depends on the scientific goal, not just convenience.
      \end{itemize}
  \end{columns}
\end{frame}

\begin{frame}[t]{Ward Linkage: The Objective Function}
  \begin{columns}[T,totalwidth=\textwidth]
    \column{0.54\textwidth}
      \footnotesize
      \begin{block}{Within-cluster error}
        \[
          \ESS(C)
          =
          \sum_{i\in C}\|x_i-\bar x_C\|_2^2
        \]
        \[
          \bar x_C = \frac{1}{|C|}\sum_{i\in C}x_i
        \]
      \end{block}

      Ward linkage merges the pair that yields the \textbf{smallest increase} in total within-cluster sum of squares:
      \[
        \Delta(A,B)
        =
        \ESS(A\cup B)-\ESS(A)-\ESS(B).
      \]
      For Euclidean distance, this simplifies to
      \[
        \Delta(A,B)
        =
        \frac{|A||B|}{|A|+|B|}
        \|\bar x_A-\bar x_B\|_2^2.
      \]
      Small centroid gaps and balanced sizes make a merge attractive.
    \column{0.40\textwidth}
      \begin{block}{Interpretation}
        \begin{itemize}
          \item Ward is a greedy least-squares method.
          \item It often prefers compact, roughly spherical clusters.
          \item This makes it feel like a hierarchical cousin of $k$-means.
        \end{itemize}
      \end{block}
  \end{columns}
\end{frame}

\begin{frame}[t]{Ward Linkage: Visual Intuition}
  \begin{columns}[T,totalwidth=\textwidth]
    \column{0.58\textwidth}
      \begin{center}
        \includegraphics[
          width=\linewidth,
          height=0.60\textheight,
          keepaspectratio
        ]{ward_linkage_ess.png}
      \end{center}
    \column{0.36\textwidth}
      \small
      \begin{itemize}
        \item First compute the error of the merged cluster.
        \item Then subtract the errors of the two original clusters.
        \item The best merge is the one with the \textbf{smallest increase} in total within-cluster variation.
      \end{itemize}
  \end{columns}
\end{frame}

\begin{frame}[t]{Different Linkages Can Tell Very Different Stories}
  \small
  \begin{center}
    \includegraphics[
      width=0.92\linewidth,
      height=0.58\textheight,
      keepaspectratio
    ]{linkage_impact_examples.png}
  \end{center}
  \begin{itemize}
    \item \textbf{Single linkage} can recover non-convex shapes, but it is prone to \emph{chaining}.
    \item \textbf{Complete} and \textbf{average linkage} usually create more balanced dendrograms.
    \item \textbf{Ward linkage} often behaves like a hierarchical cousin of $k$-means: it likes compact groups.
  \end{itemize}
\end{frame}

\begin{frame}[t]{Cutting the Dendrogram Gives a Clustering}
  \small
  \begin{columns}[T,totalwidth=\textwidth]
    \column{0.43\textwidth}
      \small
      Choose a height $h$ and remove every merge above that level. The remaining connected components define the clusters:
      \[
        \mathcal C_h
        =
        \{\text{connected components below height } h\}.
      \]

      \begin{itemize}
        \item Higher cut $\Rightarrow$ fewer clusters.
        \item Lower cut $\Rightarrow$ more clusters.
        \item In practice, we look for a large vertical gap before making the cut.
      \end{itemize}
      \emph{$K$-means chooses $K$ first; hierarchical clustering builds the whole tree first and chooses $K$ by the cut.}
    \column{0.53\textwidth}
      \begin{center}
        \includegraphics[
          width=\linewidth,
          height=0.50\textheight,
          keepaspectratio
        ]{dendrogram_cut_levels.png}
      \end{center}
  \end{columns}
\end{frame}

\begin{frame}[t]{Important Caveats from ISLP}
  \begin{columns}[T,totalwidth=\textwidth]
    \column{0.47\textwidth}
      \begin{block}{Nested clusters are an assumption}
        Cutting a dendrogram at height $h_1$ and then at a lower height $h_2$ always produces nested groups.
      \end{block}

      \begin{itemize}
        \item That nested structure is a property of the algorithm, not proof that the data are truly hierarchical.
        \item In some data sets, the best 2-group split and best 3-group split may not be consistent with one another.
      \end{itemize}
    \column{0.47\textwidth}
      \begin{block}{Practical consequence}
        Hierarchical clustering can sometimes give worse results than methods like $k$-means for a fixed number of clusters.
      \end{block}

      \begin{itemize}
        \item Greedy early merges cannot be repaired later.
        \item Compare several linkage and scaling choices rather than trusting a single tree blindly.
      \end{itemize}
  \end{columns}
\end{frame}

\begin{frame}[t]{A Tempered Approach to Interpreting Clusters}
  \begin{itemize}
    \item Small choices can have big consequences: scaling, distance measure, and linkage can all change the result.
    \item There is rarely a single ``correct'' clustering waiting to be discovered automatically.
    \item A good workflow is to compare several reasonable specifications and look for patterns that recur.
    \item If the same clusters appear across subsets of the data, that gives more confidence that the structure is robust.
    \item Treat clustering as an \textbf{exploratory tool}: a starting point for scientific interpretation, not the last word.
  \end{itemize}

  \vspace{0.6em}
  \begin{block}{Practical habit}
    Use clustering to generate hypotheses and summaries, then validate those ideas with domain knowledge or new data whenever possible.
  \end{block}
\end{frame}

\begin{frame}[t]{Strengths and Limitations}
  \begin{columns}[T,totalwidth=\textwidth]
    \column{0.47\textwidth}
      \begin{block}{Strengths}
        \begin{itemize}
          \item No need to choose one partition in advance
          \item Works with many dissimilarity measures, not just Euclidean distance
          \item Dendrograms are often easy to explain to humans
          \item Useful for exploratory analysis and small-to-medium data sets
        \end{itemize}
      \end{block}
    \column{0.47\textwidth}
      \begin{block}{Limitations}
        \begin{itemize}
          \item Distances can change dramatically if variables are not scaled
          \item Early greedy merges cannot be undone later
          \item Results depend strongly on the choice of linkage and distance metric
          \item Pairwise distances require about $O(n^2)$ memory
        \end{itemize}
      \end{block}
  \end{columns}

  \vspace{0.2em}
  \[
    \text{memory} \approx O(n^2),
    \qquad
    \text{time typically ranges from } O(n^2\log n) \text{ to } O(n^3).
  \]
\end{frame}

\begin{frame}[t]{Takeaways}
  \begin{itemize}
    \item Hierarchical clustering builds a \textbf{tree of nested clusters}, not just one partition.
    \item Agglomerative clustering is a greedy bottom-up procedure that repeatedly merges the closest clusters.
    \item The \textbf{linkage function} is the main mathematical choice because it defines what ``closest'' means for clusters.
    \item Distances, scaling, and linkage choices all shape the final dendrogram, so interpretation should be cautious.
    \item Ward linkage minimizes the increase in within-cluster sum of squares, linking hierarchical clustering to least-squares ideas.
    \item A dendrogram becomes a concrete clustering only after we choose a \textbf{cut height}.
  \end{itemize}
\end{frame}

\end{document}
